\subsection{能耗、充电与时变碳排放核算}

在给定行驶速度下，电动车逐弧电耗与燃油车逐弧油耗写为距离和载重的函数：
\begin{align}
e^E_{ij,k}&=(\alpha_{0k}^E+\alpha_{1k}^E L_{ir})d_{ij}, && k\in K^E,\label{eq:ev-energy}\\
f^F_{ij,k}&=(\alpha_{0k}^F+\alpha_{1k}^F L_{ir})d_{ij}, && k\in K^F.\label{eq:cv-fuel}
\end{align}
其中，最终稿仍需逐式核对系数的单位、实现值和文献来源\cite{ref:goeke-2015}。

设充电会话$q$发生于站点$s(q)$，功率为$P_{s(q)}$，其与电网时段$t$的重叠时长为
\begin{equation}
g_{qt}=\left[\min\{h_q+\Delta_q,\bar T_t\}-\max\{h_q,\underline T_t\}\right]_+,
\label{eq:charge-overlap}
\end{equation}
则时段充电量为$y_{qt}=P_{s(q)}g_{qt}/60$，并满足$\sum_{t\in T}y_{qt}=Y_q$。部分充电和公共站容量的处理参照相关电动车路径研究\cite{ref:froger-2022,ref:keskin-2016}。

实体车排班模式$p$的燃油车直接排放和电动车充电侧间接排放分别为
\begin{align}
E_{kp}^{F}&=\lambda^F\sum_{r\in\mathcal R_{kp}}\sum_{(i,j)\in r}f^F_{ij,k},\label{eq:cv-emission}\\
E_{kp}^{E}&=\sum_{q\in\mathcal Q_{kp}}\sum_{t\in T}\gamma_ty_{qt}.\label{eq:ev-emission}
\end{align}
系统全口径排放为
\begin{equation}
E^{\mathrm{sys}}=\sum_{k\in K}\sum_{p\in\Omega_k}(E_{kp}^{F}+E_{kp}^{E})z_{kp}.
\label{eq:system-emission}
\end{equation}
式\eqref{eq:system-emission}是后文各决策层共同使用的排放尺：移动充电开始时刻改变$y_{qt}$的时段分布，改变车型指派改变直接与间接排放的构成，动态重规划则改变车辆可用时刻和充电窗口。

\subsection{目标函数与约束}

运营成本由实体车固定成本、里程成本、燃油成本和充电电费构成：
\begin{equation}
C^{\mathrm{op}}=\sum_{k\in K}c_k^{\mathrm{fix}}u_k+
\sum_{k\in K}\sum_{p\in\Omega_k}\left(c_k^{\mathrm{km}}D_{kp}+p^FF_{kp}+
\sum_{q\in\mathcal Q_{kp}}\sum_{t\in T}p^e_{s(q),t}y_{qt}\right)z_{kp}.
\label{eq:operation-cost}
\end{equation}
固定成本按启用实体车计取：
\begin{equation}
u_k=\sum_{p\in\Omega_k}z_{kp},\qquad \sum_{p\in\Omega_k}z_{kp}\leq 1.
\label{eq:vehicle-use}
\end{equation}
采用单位碳价把排放纳入统一货币目标：
\begin{equation}
\min Z=C^{\mathrm{op}}+p^{\mathrm{car}}E^{\mathrm{sys}}.
\label{eq:objective}
\end{equation}

每个客户由一个实体车排班模式恰好覆盖：
\begin{equation}
\sum_{k\in K}\sum_{p\in\Omega_k}A_{ikp}z_{kp}=1,\qquad i\in N.
\label{eq:coverage}
\end{equation}
模式内部满足流守恒、载重与硬时间窗约束：
\begin{align}
\sum_jx_{ijr}&=\sum_jx_{jir}=a_{ir}, && i\in N,\label{eq:flow}\\
0\leq L_{ir}&\leq Q_k,\qquad L_{jr}\leq L_{ir}-q_j+M(1-x_{ijr}),\label{eq:load}\\
e_i a_{ir}&\leq T_{ir}\leq l_i+M(1-a_{ir}),\nonumber\\
T_{jr}&\geq T_{ir}+s_i+t_{ij}-M(1-x_{ijr}).\label{eq:time-window}
\end{align}
电动车电量沿弧传播，并仅在允许地点补能：
\begin{equation}
b_k^{\min}\leq b_{ir}\leq B_k,\qquad
b_{jr}\leq b_{ir}+Y_{ir}-e^E_{ij,k}+M(1-x_{ijr}),\quad k\in K^E.
\label{eq:soc}
\end{equation}
同一实体车相邻趟$r\rightarrow r'$满足时间与电量连续：
\begin{align}
T_{d^+_{r'}r'}&\geq T_{d^-_rr}+\sum_{q\in\mathcal Q_{kp}^{rr'}}\Delta_q,\label{eq:trip-time}\\
b_{d^+_{r'}r'}&=b_{d^-_rr}+\sum_{q\in\mathcal Q_{kp}^{rr'}}Y_q.\label{eq:trip-energy}
\end{align}
公共站容量约束为
\begin{equation}
\sum_{k\in K}\sum_{p\in\Omega_k}O_{skpt}z_{kp}\leq C_s,\qquad s\in S,\ t\in T.
\label{eq:station-capacity}
\end{equation}
式\eqref{eq:station-capacity}不施加于车场充电。

设第$m$次重规划时刻为$\tau_m$，待服务客户集合按事件更新：
\begin{equation}
N^{\tau_m}=\left(N^{\tau_{m-1}}\setminus N_{\mathrm{served}}^{\tau_m}\setminus N_{\mathrm{cancel}}^{\tau_m}\right)\cup N_{\mathrm{new}}^{\tau_m}.
\label{eq:dynamic-set}
\end{equation}
已经发生的成本与排放分别累计为
\begin{equation}
\bar C^{\tau_m}=\bar C^{\tau_{m-1}}+\Delta C^{\tau_{m-1}},\qquad
\bar E^{\tau_m}=\bar E^{\tau_{m-1}}+\Delta E^{\tau_{m-1}}.
\label{eq:dynamic-account}
\end{equation}
新一轮优化仅计算剩余问题，已执行前缀作为不可撤销约束。

\section{算法设计}

\subsection{算法框架}

GCI-DMM-VRP同时包含路径、车型、实体车多趟、充电和动态状态继承。本文采用成熟混合遗传搜索作为搜索底座，在其上组织一条面向本问题的串行改进链\cite{ref:vidal-2012,ref:wouda-2024}。底座本身不作为算法创新；问题特化组件的最终贡献由被接受的增量式消融决定。

框架坚持“唯一一本账”：候选是否被接受、能否进入种群、能否存活以及能否被选为父代，均由同一完整评价器给出的完整成本与完整可行性决定。能源消耗、充电、跨趟时间---电量衔接、时变电价和时变碳强度进入同一本账。候选完成失败保留失败原因，不静默冒充初始解。

\begin{figure}[htbp]
  \centering
  \begin{tikzpicture}[
    node distance=5mm and 7mm,
    box/.style={draw=black!55,rounded corners,fill=black!4,align=center,minimum height=8mm,text width=28mm,font=\small},
    arrow/.style={-{Latex[length=2mm]},thick,draw=black!65}
  ]
    \node[box] (snapshot) {静态输入或动态状态快照};
    \node[box,right=of snapshot] (initial) {完整可行起点};
    \node[box,right=of initial] (hgs) {选择、交叉、教育与多样性管理};
    \node[box,below=of hgs] (complete) {实体车多趟与充电补全};
    \node[box,left=of complete] (evaluate) {完整成本与完整可行性评价};
    \node[box,left=of evaluate] (population) {按完整结果入群、生存与父代选择};
    \node[box,below=of evaluate] (audit) {独立验解并保存最好解};
    \draw[arrow] (snapshot)--(initial);
    \draw[arrow] (initial)--(hgs);
    \draw[arrow] (hgs)--(complete);
    \draw[arrow] (complete)--(evaluate);
    \draw[arrow] (evaluate)--(population);
    \draw[arrow] (population.west) -- ++(-6mm,0) |- (hgs.north);
    \draw[arrow] (evaluate)--(audit);
  \end{tikzpicture}
  \caption{算法总体流程。流程图给出已确定的科学合同；\TBD{算法定稿后回填最终组件清单与顺序}。}
  \label{fig:algorithm-flow}
\end{figure}

\subsection{算法步骤}

\textbf{步骤一：构造状态快照与可行起点。} 静态阶段读取订单、车场、有限实体车队和碳强度序列；动态阶段读取式\eqref{eq:dynamic-set}--\eqref{eq:dynamic-account}定义的剩余客户与车辆状态。初始个体先经完整模型检查。

\textbf{步骤二：执行混合遗传搜索主循环。} 算法维护可行与不可行子种群，通过选择、交叉、教育与多样性管理更新候选，罚系数按可行解比例调整。停止条件和预算由本项目收敛记录标定，不直接照搬其他算法的迭代数。

\textbf{步骤三：施加问题特化动作。} 便宜代理只提名有界候选，完整真值逐项评价并决定是否接受。\TBD{算法定稿后列出进入论文的特化动作与出处}。

\textbf{步骤四：完成实体车排班与充电。} 将配送趟分配给具体实体车，按时间顺序衔接多趟并补齐电量、充电会话、分时电价和碳排放。违反时间窗、容量、电量、公共站容量或实体车上限的候选保留失败原因。

\textbf{步骤五：入群、生存与父代选择。} 入群按完整可行性分池，生存选择和父代选择使用同一完整账；全局最好解必须完整可行。

\textbf{步骤六：独立验解与动态调用。} 进入种群的候选重新执行独立检查，只有完整可行且严格改善的候选更新最好解。事件触发后冻结已经执行的路径前缀，更新车辆状态，再求解剩余问题。

\section{数值实验设计与算法有效性}

\subsection{算例、参数与代表解}

当前算例套件保留“China81”历史名称，但权威清单含82个目录条目：成渝27个、京津冀28个、珠三角27个；额外条目是一个京津冀50客户候选。代表实例为京津冀50客户、2车场、97个公共站节点的候选，客户责任按编号等分为两家企业各25个。动态事件流包含10个新增客户事件，接单窗为08:00--10:00，共5个触发批；分时电价和碳强度的信息分辨率为1小时。

\begin{table}[htbp]
  \centering
  \caption{算例规模与结构健康}
  \label{tab:instances}
  \scriptsize
  \begin{tabular}{p{0.30\linewidth}rrrr}
    \toprule
    指标 & 成渝 & 京津冀 & 珠三角 & 全套件\\
    \midrule
    目录条目数（个） & 27 & 28 & 27 & 82\\
    客户规模档（个） & \multicolumn{4}{c}{10、15、20、25、50、75、100、150、200}\\
    每档复本数（个） & 3 & 3（50客户另有候选） & 3 & 3（另有候选）\\
    完成客户数/总客户数 & \TBD{套件正式搜索} & \TBD{套件正式搜索} & \TBD{套件正式搜索} & \TBD{套件正式搜索}\\
    完成需求量/总需求量 & \TBD{套件正式搜索} & \TBD{套件正式搜索} & \TBD{套件正式搜索} & \TBD{套件正式搜索}\\
    每例车场数 & \TBD{汇总} & \TBD{汇总} & \TBD{汇总} & \TBD{汇总}\\
    每例公共站数 & \TBD{汇总} & \TBD{汇总} & \TBD{汇总} & \TBD{汇总}\\
    时间窗宽度 & \TBD{汇总} & \TBD{汇总} & \TBD{汇总} & \TBD{汇总}\\
    上午/下午订单数 & \TBD{汇总} & \TBD{汇总} & \TBD{汇总} & \TBD{汇总}\\
    可争夺结构检查 & \TBD{当前套件复核} & \TBD{当前套件复核} & \TBD{当前套件复核} & \TBD{当前套件复核}\\
    两班次完整服务见证 & \TBD{当前套件复核} & \TBD{当前套件复核} & \TBD{当前套件复核} & \TBD{当前套件复核}\\
    混合车队“两侧”条件 & \TBD{当前套件复核} & \TBD{当前套件复核} & \TBD{当前套件复核} & \TBD{当前套件复核}\\
    \bottomrule
  \end{tabular}
  \tabnote{注：前3行来自当前算例清单。其余行保留已批准表壳，但当前获准数据源没有相应汇总，故不移用旧套件数字。82个目录条目是构造套件条目数，不是企业实测样本数。}
\end{table}

\begin{table}[htbp]
\centering
\caption{关键模型与算法参数及证据身份}
\label{tab:parameters}
\scriptsize
\resizebox{\textwidth}{!}{%
\begin{tabular}{llll}
\toprule
参数 & 数值与单位 & 适用层 & 证据身份/当前处置\\
\midrule
随机种子 & 8 & 代表解复现 & 保存解身份，算法定稿前读数\\
交叉模式 & fast\_only & 交叉 & 当前运行配置；出处待补\\
教育深度 & 1层 & 局部改进 & 标定依据待补，不写成最终值\\
停滞耐心 & 500次无改善 & 停止/重启 & 当前运行配置；待正式标定\\
初始/最小/最大罚系数 & 100.0/0.1/100000.0 & 可行性罚项 & 开源上游默认值\\
罚系数降低/提高倍数 & 0.32/1.34 & 可行性罚项 & 开源上游默认值\\
更新间隔/目标可行比例/容差 & 50/0.43/0.05 & 可行性罚项 & 开源上游默认值\\
修复概率/增强倍数 & 0.8/12 & 不可行解修复 & 开源上游配置\\
最小/最大种群 & 25/65个个体 & 种群管理 & 最大值由25+40得到\\
每代生成数/精英数/邻近数 & 40/4/5 & 种群管理 & 开源上游配置\\
多样性下界/上界 & 0.1/0.5 & 种群多样性 & 开源上游配置\\
锦标赛规模 & 2个个体 & 父代选择 & 当前配置；出处待补\\
充电量/时刻策略 & just\_enough/cost\_plus\_carbon & 充电 & 当前源解实际配置\\
碳价 & 0.07502元/kgCO$_2$e & 总目标 & 全国碳市场价格情景；原始页码待补\\
车场分时电价 & 谷0.56328575、平0.83644275、峰1.14862175元/kWh & 车场充电 & 北京2025-02-12数据版本\\
公共站服务费 & 0.4元/kWh & 公共站充电 & 构造情景，不写成市场实测\\
柴油价 & 7.48元/L & 燃油运营 & 北京地方零售表数据版本\\
逐时碳强度 & 0.15410000--0.64390000 kgCO$_2$e/kWh & 充电排放 & 北京2025-02-12，24个独立小时值\\
燃油/电动车非能源里程成本 & 0.7800/0.9145元/km & 车型与运营成本 & 电动车值含0.2445元/km电池折旧\\
燃油/电动车日固定成本 & 170.00/220.00元/(辆$\cdot$日) & 车辆启用 & 电动车含50元日固定溢价\\
车场单车额定充电功率 & 60.0 kW & 充电可行性 & 中国官方规划取值\\
车辆容积容量 & 7.2 m$^3$ & 路线可行性 & 当前数据集给定；出处待补\\
两车场车队上限 & 3CV+3EV；7CV+7EV & 车型与车辆启用 & 两班次物理见证派生\\
\bottomrule
\end{tabular}
}
\tabnote{注：算法参数只描述表\ref{tab:final-solution}源解的真实运行身份；标为“待补”的来源不因进入本表而升级为正式参数证据。}
\end{table}

\begin{figure}[htbp]
  \centering
  \begin{minipage}{0.48\linewidth}
  \centering
  \begin{tikzpicture}
  \begin{axis}[
    width=\linewidth,height=58mm,
    xlabel={经度},ylabel={纬度},
    tick label style={font=\scriptsize},label style={font=\small},
    legend style={font=\scriptsize,at={(0.02,0.98)},anchor=north west,draw=none,fill=white,fill opacity=.8},
    grid=both,grid style={black!8}]
    \addplot[only marks,mark=+,mark size=1.2pt,gray!45]
      table[x=lon,y=lat,col sep=comma,restrict expr to domain={\thisrow{kind}}{2:2}]{data/figure2_nodes_plot.csv};
    \addlegendentry{公共站}
    \addplot[only marks,mark=o,mark size=1.8pt,blue!65]
      table[x=lon,y=lat,col sep=comma,restrict expr to domain={\thisrow{kind}}{0:0},restrict expr to domain={\thisrow{enterprise}}{0:0}]{data/figure2_nodes_plot.csv};
    \addlegendentry{企业甲客户}
    \addplot[only marks,mark=triangle*,mark size=1.7pt,orange!85!black]
      table[x=lon,y=lat,col sep=comma,restrict expr to domain={\thisrow{kind}}{0:0},restrict expr to domain={\thisrow{enterprise}}{1:1}]{data/figure2_nodes_plot.csv};
    \addlegendentry{企业乙客户}
    \addplot[only marks,mark=square*,mark size=3pt,black]
      table[x=lon,y=lat,col sep=comma,restrict expr to domain={\thisrow{kind}}{1:1}]{data/figure2_nodes_plot.csv};
    \addlegendentry{车场}
  \end{axis}
  \end{tikzpicture}
  \\(a) 节点分布
  \end{minipage}\hfill
  \begin{minipage}{0.48\linewidth}
  \centering
  \begin{tikzpicture}
  \begin{axis}[
    width=\linewidth,height=58mm,
    xlabel={经度},ylabel={纬度},
    tick label style={font=\scriptsize},label style={font=\small},
    grid=both,grid style={black!8},unbounded coords=jump,
    legend style={font=\scriptsize,at={(0.02,0.98)},anchor=north west,draw=none,fill=white,fill opacity=.8}]
    \addplot[thin,blue!60] table[x=lon,y=lat,col sep=comma]{data/figure2_routes_a.csv};
    \addlegendentry{企业甲所属车}
    \addplot[thin,orange!80!black] table[x=lon,y=lat,col sep=comma]{data/figure2_routes_b.csv};
    \addlegendentry{企业乙所属车}
    \addplot[only marks,mark=o,mark size=1.3pt,black!65]
      table[x=lon,y=lat,col sep=comma,restrict expr to domain={\thisrow{kind}}{0:0}]{data/figure2_nodes_plot.csv};
    \addplot[only marks,mark=square*,mark size=3pt,black]
      table[x=lon,y=lat,col sep=comma,restrict expr to domain={\thisrow{kind}}{1:1}]{data/figure2_nodes_plot.csv};
  \end{axis}
  \end{tikzpicture}
  \\(b) 路线及车型
  \end{minipage}
  \caption{代表实例的节点分布与路线。节点表含2个车场、50个客户和97个公共站节点；保存解包含16趟路线，均由燃油车执行。路线未作手工平滑或美化。}
  \label{fig:nodes-routes}
\end{figure}

\begin{table}[htbp]
  \centering
  \caption{代表实例解的成本、排放与资源分解}
  \label{tab:final-solution}
  \small
  \begin{tabular}{llr}
    \toprule
    类别 & 指标（单位） & 保存解数值\\
    \midrule
    服务 & 完成客户数/总客户数（个） & 50/50\\
    服务 & 完成需求量/总需求量（kg） & 13264/13264\\
    资源 & 启用燃油车/电动车（辆） & 8/0\\
    资源 & 启用实体车总数（辆） & 8\\
    资源 & 配送趟次（趟） & 16\\
    里程 & 总里程（km） & 935.530\\
    成本 & 车辆固定成本（元） & 1360.000\\
    成本 & 非能源里程成本（元） & 729.713\\
    成本 & 燃油成本（元） & 919.139\\
    成本 & 电费（元） & 0.000\\
    成本 & 碳价折算成本（元） & 24.354\\
    成本 & 单目标总账（元） & 3033.206\\
    能源 & 燃油消耗（L） & 122.879\\
    能源 & 电能消耗（kWh） & 0.000\\
    排放 & 燃油车直接排放（kgCO$_2$e） & 324.636\\
    排放 & 电动车充电排放（kgCO$_2$e） & 0.000\\
    可行性 & 完整评价可行 & 是\\
    可行性 & 硬约束违规数（项） & 0\\
    \bottomrule
  \end{tabular}
  \tabnote{注：保存解来自协同效率批的联合运行，公平约束关闭。该解用于展示一份可执行方案如何闭合服务、资源、成本与排放，不承担算法优越性证明。总里程由8辆实体车的逐车里程求和得到，与联合解分解字段一致。}
\end{table}

\begin{table}[htbp]
  \centering
  \caption{代表解实体车辆的多趟安排}
  \label{tab:vehicle-schedules}
  \scriptsize
  \resizebox{\textwidth}{!}{%
  \begin{tabular}{llllrr}
    \toprule
    实体车 & 所属车场 & 车型 & 趟数 & 服务客户数 & 里程（km）\\
    \midrule
    CV\_A\_1 & 车场A & 燃油车 & 2 & 6 & 151.083\\
    CV\_A\_2 & 车场A & 燃油车 & 1 & 2 & 83.613\\
    CV\_B\_1 & 车场B & 燃油车 & 1 & 4 & 56.999\\
    CV\_B\_2 & 车场B & 燃油车 & 3 & 10 & 151.064\\
    CV\_B\_3 & 车场B & 燃油车 & 3 & 8 & 162.811\\
    CV\_B\_4 & 车场B & 燃油车 & 2 & 6 & 93.774\\
    CV\_B\_5 & 车场B & 燃油车 & 2 & 7 & 119.043\\
    CV\_B\_6 & 车场B & 燃油车 & 2 & 7 & 117.142\\
    \bottomrule
  \end{tabular}
  }
  \tabnote{注：逐车客户序列已经保存，但当前展品数据没有逐趟发车/返回时刻、油耗、直接排放与运营成本字段；这些项目均为\TBD{重新完整评价保存解后提取}。}
\end{table}
